% !TeX root = ../main.tex
\chapter{Object-Oriented Design Patterns}
\label{ch:design-patterns}

Architectural patterns describe how the system is organised. The patterns in
this chapter are smaller: they describe how particular problems inside the
subsystems are solved. The names, intents and consequences are those of Gamma
and colleagues, listed in Chapter~\ref{ch:references}. Five are adopted and five
were considered and not adopted. In two cases the pattern is not a choice we
made but one the vendor API requires us to meet. We say so where that applies,
rather than present it as a decision.

\section{Classification}

Each pattern has a purpose, which is what it is for, and a scope, which says
whether it works through inheritance between classes or through composition
between objects. Setting the adopted set out this way shows where the design's
difficulty actually lies.

\begin{longtable}{L{0.17\textwidth}L{0.23\textwidth}L{0.52\textwidth}}
	\caption{The adopted patterns by purpose and scope}\label{tab:classification}\\
	\toprule
	\textbf{Pattern} & \textbf{Purpose and scope} & \textbf{Problem it solves here} \\
	\midrule
	\endfirsthead
	\toprule
	\textbf{Pattern} & \textbf{Purpose and scope} & \textbf{Problem it solves here} \\
	\midrule
	\endhead

	\textsc{Composite} & Object structural & Behaviour constructs nest, and the
	code that walks them should not care whether it holds a leaf or a
	branch. \\
	\addlinespace
	\textsc{Adapter} & Object structural & Two vendor libraries speak their own
	vocabularies, and neither is ours. \\
	\addlinespace
	\textsc{Visitor} & Object behavioural & An operation over the construct tree
	should be added without changing the construct classes. \\
	\addlinespace
	\textsc{Strategy} & Object behavioural & Two target languages, one
	generation algorithm. \\
	\addlinespace
	\textsc{Command} & Object behavioural & A change must be describable before
	it is performed. \\
	\bottomrule
\end{longtable}

Two things follow from the table. Every adopted pattern has object scope rather
than class scope, so the design is held together by composition rather than by
inheritance hierarchies. That is consistent with the pipeline of
Chapter~\ref{ch:patterns}, where stages are assembled rather than subclassed.

More telling, no creational pattern is adopted at all. Creational patterns solve
the problem of deciding which objects to make and how to make them. That is not
where this system is difficult. The set of things the tool creates is fixed in
advance by \textbf{C-11} on the model side and by the target language's grammar
on the code side, so there is nothing to defer or configure. The difficulty is
in traversing a structure and translating it faithfully, which is why the
adopted set is structural and behavioural throughout.

\section{Composite and Visitor: The Pair That Carries the Behaviour}
\label{sec:composite-visitor}

These two are treated together because neither is useful here without the other.

\subsection{Composite}

Sequencing inside a GENESYS behaviour diagram is carried by the behaviour
constructs Function, Loop, Loop Exit, Iterate, Select, Parallel, Replicate and
Exit, which is recorded as \textbf{C-11}. These nest: a Loop contains the
constructs that form its body, a Select contains its branches, and either may
contain further loops and selects.

That is the \textsc{Composite} intent exactly, which is to compose objects into
tree structures representing part-whole hierarchies and let clients treat
individual objects and compositions uniformly. The abstract syntax in
\code{Pentateuch.Absyn} is a tree of this shape, and it is what lets
\req{CA-2.4} be expressed as a single construct rather than an unrolled sequence
of nodes.

The pattern's stated cost is that it can make a design overly general, and in
particular that restrictions such as type checking become harder to express. We
accept that cost, because the restriction we care about is not structural but
semantic: a construct that cannot be translated is reported under \req{CA-8} and
\req{AS-3} rather than prevented from being built.

\subsection{Visitor}

Having a composite tree, something must operate on it, and the operations differ
by direction. Generation turns each construct into a control flow structure in
the target language. Extraction produces constructs from a control flow graph.
Both dispatch on the construct type, and both would otherwise be a chain of type
tests spread across the code that walks the tree.

\textsc{Visitor} separates the operation from the structure it operates on, so a
new operation is a new visitor rather than a change to every node class. This is
not a free choice. The GENESYS structure traversal interface is itself a visitor
interface, and the set of constructs in \textbf{C-11} is precisely the set of
methods it exposes. We implement the pattern because the vendor API presents it.

The fit is closer than convenience. The pattern's stated condition for use is
that the classes defining the object structure rarely change while new
operations over the whole structure are often needed. The construct set is
defined by the vendor and fixed for any given GENESYS release, and each new
translation direction is a new operation over it. That is the condition, almost
word for word.

That matters because the pattern's main drawback is the reverse case: adding a
new element class is hard, since every visitor must gain a method for it. If
GENESYS introduced a ninth behaviour construct, every visitor in the tool would
have to change. \textbf{C-11} is what makes that drawback acceptable rather than
dangerous, since the set is not ours to extend and does not move.

This arrangement gives us a property we rely on. Because the visitor interface
enumerates the construct types, a construct the tool does not handle cannot be
silently ignored: the method exists and must do something. That is what makes
\req{CA-8} and \req{AS-3}, reporting a construct rather than mistranslating it,
implementable with confidence rather than by hoping the type tests were
exhaustive.

\section{Adapter: Two Vendor APIs, One Vocabulary}

Neither the Client API nor Roslyn speaks in terms of our abstract syntax, and we
control neither. \code{Pentateuch.Genesys} wraps the Client API and
\code{Pentateuch.Translator} wraps Roslyn, so that each presents its data in our
terms.

This is the \textsc{Adapter} pattern: converting one interface into the one a
client expects. The incompatibility is real, since \textbf{C-1} means the two
libraries cannot be referenced from the same project. The adapters are therefore
also where the solution splits, so the subsystem boundaries of
Chapter~\ref{ch:subsystems} and the adapter boundaries coincide.

\paragraph{Object adapter, not class adapter.}
We use the object form, which holds a reference to the adapted type instead of
inheriting from it. A class adapter cannot cover a type together with all of its
subtypes, and the GENESYS entity classes form hierarchies that \textbf{C-7} says
we must discover at run time. One object adapter can serve the whole family.

\paragraph{Why not Facade.}
A facade simplifies access to a subsystem but keeps its vocabulary. An adapter
changes the vocabulary. Our wrappers do simplify, but that is not why they
exist: everything above them is written against the abstract syntax and must not
know whether a model came from the Client API or from a file. That is interface
conversion, so \textsc{Adapter} is the right name.

A further benefit follows. Because the layers above use our vocabulary, the
interchange-file fallback of \textbf{R-1} means replacing one adapter and
nothing else, as described in Section~\ref{sec:layered}.

\section{Strategy: Two Target Languages, One Pipeline}

\req{AC-1} requires the tool to generate C\# or Visual Basic. The difference
between them is confined to how each construct is written out, and nothing
upstream of generation changes.

\textsc{Strategy} defines a family of algorithms, encapsulates each, and makes
them interchangeable, so the Code Generator holds one algorithm and two
strategies rather than two algorithms. Combined with the canonical form of
Section~\ref{sec:canonical}, this is what makes a second language additive
and interchangeable.

The pattern's stated costs apply only weakly here. One is that clients must know
about all the strategies in order to choose between them, which in our case
means the user names a target language and nothing else has to. The other is an
increase in the number of objects, which against a code generator that already
holds a visitor per direction is not a meaningful addition.

\section{Command: Proposing a Change Before Making It}
\label{sec:command}

\req{CM-3} requires the tool to present what it proposes to change and to write
nothing until the user accepts. \req{NFR-3} requires that a cancelled or failed
operation leave nothing partially written, on either side.

Both are satisfied by representing each intended modification as an object
rather than performing it immediately. \textsc{Command} encapsulates a request
in exactly this way, and the pattern's stated uses include specifying and
executing requests at different times, logging changes, and structuring a system
around high-level operations built from primitive ones. All three describe what
the review step does.

The change set shown at review is the collected list of such objects, which is
what allows it to be displayed before any of it is applied. Acceptance executes
the list; cancellation discards it. Because nothing has been written when the
decision is made, the property \req{NFR-3} asks for is structural rather than
something that has to be undone afterwards. Test V-9 checks it by comparing the
repository byte for byte against its prior state.

Two further consequences come for free. The pattern notes that commands can be
assembled into composite commands, which is what a change set is: a single
acceptance applies a list of individual modifications as one unit. And the same
objects carry enough information to describe what was done, which is what
\req{NFR-7} asks a log to contain.

\section{Patterns Considered and Not Adopted}

\begin{longtable}{L{0.19\textwidth}L{0.73\textwidth}}
	\caption{Design patterns considered and not adopted}\label{tab:rejected-patterns}\\
	\toprule
	\textbf{Pattern} & \textbf{Why not} \\
	\midrule
	\endfirsthead
	\toprule
	\textbf{Pattern} & \textbf{Why not} \\
	\midrule
	\endhead

	\textsc{Interpreter} & The closest call of the five. The pattern's own
	description of its structure is ours: the abstract syntax tree is a
	\textsc{Composite}, and a \textsc{Visitor} holds the behaviour for each
	node in one class. We have both. What we do not have is its purpose, which
	is to interpret sentences in a language. Nothing evaluates our tree; it is
	translated into another representation and discarded. Adopting the name
	would describe the data structure correctly and the behaviour not at
	all. \\
	\addlinespace
	\textsc{Template Method} & Names something real. All four translations
	share a skeleton: connect, read, abstract, map, review, write, with the
	middle varying by direction. \textsc{Template Method} is the way to express
	that through inheritance, with a base class fixing the order and subclasses
	filling in steps. We express the same commonality by composition instead,
	as shared filters in one pipeline, which is what
	Chapter~\ref{ch:patterns} describes. The pattern is not wrong about the
	structure; the architecture already carries it. \\
	\addlinespace
	\textsc{Singleton} & Superficially attractive for the repository
	connection, since \textbf{C-5} means one project is open at a time and the
	connection is expensive to establish. Rejected because a global instance
	would make the connection impossible to substitute in a test, and
	\req{NFR-6} exists specifically so that the projects not referencing
	GENESYS can be tested without it. A single instance passed explicitly gives
	the same lifetime without the coupling. \\
	\addlinespace
	\textsc{Memento} & The usual companion to \textsc{Command}, used to capture
	enough state to undo an operation. Not needed, because the design avoids
	the problem it solves. Nothing is written until the user accepts, so there
	is no committed state to roll back and no snapshot to keep. This is
	\req{NFR-3} achieved by ordering rather than by reversal. \\
	\addlinespace
	\textsc{Observer} & Would suit \req{FS-2}, the status window that follows
	each translation and reports progress and outcome. Not adopted now, because
	the minimum viable product performs one translation per invocation and
	reports at the end, so there is no ongoing state for an observer to watch.
	It would become relevant if \req{FS-2} is built. \\
	\bottomrule
\end{longtable}

\section{Naming Convention for the Implementation}

Where a class exists because of a pattern, its name ends in the pattern's name,
so that \code{ConstructVisitor}, \code{GenesysEntityAdapter} and
\code{CreateEntityCommand} are self-describing to a reader who knows the
patterns and merely verbose to one who does not. This costs nothing and makes
the chapter above checkable against the source during the build semester, which
matters because a pattern claimed in a design document and absent from the code
is worse than no claim at all.

\section{Summary}

\begin{longtable}{L{0.17\textwidth}L{0.42\textwidth}L{0.33\textwidth}}
	\caption{Design patterns adopted}\label{tab:patterns-adopted}\\
	\toprule
	\textbf{Pattern} & \textbf{Where} & \textbf{What it serves} \\
	\midrule
	\endfirsthead
	\toprule
	\textbf{Pattern} & \textbf{Where} & \textbf{What it serves} \\
	\midrule
	\endhead

	\textsc{Composite} & The abstract syntax tree of the behaviour constructs
	in \textbf{C-11} & \req{CA-2.4}, representing a loop once \\
	\addlinespace
	\textsc{Adapter} & \code{Pentateuch.Genesys} over the Client API,
	\code{Pentateuch.Translator} over Roslyn, object form in both &
	\textbf{C-1}, \textbf{C-7}, and the fallback in risk \textbf{R-1} \\
	\addlinespace
	\textsc{Visitor} & Dispatch over those constructs, following the vendor's
	own traversal interface & \req{CA-8} and \req{AS-3}, reporting rather than
	mistranslating \\
	\addlinespace
	\textsc{Strategy} & Emission rules inside the Code Generator &
	\req{AC-1}, two target languages \\
	\addlinespace
	\textsc{Command} & The change set held at the review step & \req{CM-3},
	\req{NFR-3}, and a path to \req{NFR-7} \\
	\bottomrule
\end{longtable}
